\section{Recomendaciones de política pública}

Las siguientes medidas son propuestas derivadas del experimento. No se
presentan como programas vigentes ni como decisiones ya adoptadas.

\begin{enumerate}
  \item \textbf{Crear un conjunto nacional con procedencia verificable.}
  Instituciones académicas, asociaciones de productores y servicios de
  extensión pueden acordar un esquema común para fotografías de hojas, con
  municipio a una resolución apropiada, variedad, etapa fenológica, fecha,
  dispositivo y diagnóstico de referencia. La diversidad regional debe
  medirse, no asumirse.

  \item \textbf{Estandarizar la captura sin volverla artificial.}
  Cada registro debería incluir una vista cercana del síntoma y otra de la
  planta, bajo luz disponible, además de una guía de enfoque. El protocolo
  necesita ejemplos reales de fondos complejos; llenar el dataset de fondos
  blancos mejoraría el laboratorio y empeoraría la pregunta de campo.

  \item \textbf{Establecer revisión agronómica de etiquetas.}
  Una corrección enviada desde la app no debe convertirse directamente en
  ground truth. Se requiere doble revisión para casos dudosos, registro del
  criterio usado y una categoría de desacuerdo. Las deficiencias N/P/K merecen
  prioridad porque su similitud visual fue el principal problema de ambas
  etapas.

  \item \textbf{Usar el sistema como apoyo de extensión.}
  El primer despliegue debería entregar la herramienta a equipos técnicos y a
  grupos acompañados de productores. Esto permite contrastar la predicción con
  síntomas de la parcela, corregir recomendaciones y definir cuándo la app debe
  derivar a una visita.

  \item \textbf{Monitorear por región, clase y dispositivo.}
  Accuracy global no basta. Un tablero de operación debe seguir recall de
  clases críticas, tasa de rechazo, correcciones, latencia y brechas entre
  cámaras o zonas. Los grupos con soporte pequeño deben mostrarse con su
  incertidumbre y no ocultarse bajo promedios nacionales.

  \item \textbf{Minimizar y proteger ubicación.}
  La georreferencia debe ser opcional, separada del uso básico y explicada en
  lenguaje claro. Cuando el objetivo sea vigilancia regional, puede
  almacenarse una cuadrícula o municipio en vez de coordenadas exactas. Deben
  definirse retención, acceso, eliminación y respuesta ante filtraciones.

  \item \textbf{Financiar el ciclo de retroalimentación.}
  Recoger casos difíciles sin personal que los revise crea una cola, no un
  dataset. El presupuesto debe incluir validación, curación, deduplicación,
  versionado y reentrenamiento, además del desarrollo de la aplicación.

  \item \textbf{Exigir validación de campo antes del uso masivo.}
  Una versión candidata debe aprobar criterios predefinidos de rendimiento por
  clase y entorno, prueba de teléfonos, estabilidad offline y comprensión del
  mensaje de incertidumbre. Una caída grande en potasio o en una región debe
  detener la promoción aunque el accuracy global siga alto.

  \item \textbf{Publicar una ficha de cada versión.}
  La ficha debe incluir corpus, fingerprint, clases, protocolo, métricas por
  grupo, limitaciones, hash del modelo y cambios desde la versión anterior,
  siguiendo la idea de reportes de modelo \citep{mitchell2019modelcards}.
  Esto permite saber qué produjo un diagnóstico meses después.
\end{enumerate}

Estas recomendaciones colocan el mayor esfuerzo donde el experimento encontró
la mayor fragilidad: datos con contexto, revisión y seguimiento. Comprar más
capacidad de cómputo puede facilitar el entrenamiento, pero no resuelve una
etiqueta ambigua ni una región ausente. La política útil conecta el modelo con
una red de asistencia y con reglas para corregirlo.
